% =====================================================================
%  sections/09_appendixes.tex  --  X. APPENDIXES
% =====================================================================
\clearpage
\section{APPENDIXES}
\label{sec:appendixes}

% ---------------------------------------------------------------------
\subsection{Entity-Relationship Diagram}
\label{sub:erm-diagram}

The diagram below documents the underlying database schema that powers
\projname{}, including the Kyungrinara financial dataset referenced
throughout this report. It maps the core tables together with their
primary and foreign key relationships, providing the structural
reference that the schema retrieval and SQL generation steps described
in earlier sections rely on.

\projfigure{0.95}{fig-erm-diagram}{Entity-relationship diagram of \protect\projnameplain{}'s database schema, showing the core tables and their relationships.}{fig:erm-diagram}

% ---------------------------------------------------------------------
\subsection{Text2SQL Service}
\label{sub:text2sql-service}

\code{text2sql\_service.py} implements the chat entry point that
converts a user question into an SQL-backed answer. It loads prior
chat history as context, invokes the LangGraph agent, extracts the
explanation, visualization payload, and follow-up questions from the
agent's response, persists the exchange to chat history, and falls
back to an error or timeout message if any step fails.

\begin{lstlisting}[style=python, numbers=none, caption={Text2SQL chat service (\code{text2sql\_service.py})}, label={lst:text2sql-service}]
from uuid import UUID
from typing import List
import re
import uuid

from langchain_core.messages import AIMessage, HumanMessage, BaseMessage
from sqlyst_app.core.exceptions.exception import NotFoundException
from sqlyst_app.llm.agent_graph import create_agent_graph
from sqlalchemy.orm import Session
from sqlyst_app.crud.chat_histories import ChatHistoriesCrud
from sqlyst_app.crud.chat_session import ChatSessionCrud
from sqlyst_app.schemas.chat_histories import ChatHistoriesRequest, ChatHistoriesResponse
from sqlyst_app.schemas.chat_schema import ChatResponse
from sqlyst_app.utils.response_utils import clean_response_for_display, extract_follow_up_questions_from_response, extract_json_from_response, keep_only_insights

agent = None
crud = ChatHistoriesCrud()
session_crud = ChatSessionCrud()

def get_agent():
    global agent
    if agent is None:
        agent = create_agent_graph()
    return agent


class Text2SQLService:
    """Service to generate SQL queries from natural language."""

    def _convert_chat_history_to_messages(self, chat_histories: List[ChatHistoriesResponse]) -> List:
        """Convert chat history to context messages using only user query and SQL query."""
        messages = []

        for i, history in enumerate(chat_histories):

            # Add user message
            messages.append(HumanMessage(content=history.user_message))

            # Add AI response with only SQL query context (not full explanation)
            if history.model_response:

                if isinstance(history.model_response, dict):
                    # Extract SQL query for context - check multiple possible locations
                    sql_query = ""

                    # Check direct sql_query field
                    if "sql_query" in history.model_response:
                        sql_query = history.model_response["sql_query"]


                    # Create minimal AI response with just SQL context
                    if sql_query:
                        ai_content = f"Previous query used: {sql_query}"
                        messages.append(AIMessage(content=ai_content))
                    else:
                        # Still add a generic AI message to maintain conversation flow
                        messages.append(AIMessage(content="Previous response generated"))
                else:
                    print(f"Model response is not a dict: {type(history.model_response)}")

        return messages

    def _save_chat_history(self, query: str, session_id: UUID, model_response: dict, db: Session, visual_source_id: UUID = None) -> UUID:
        """
        Save chat history with visual_source creation.
        If visual_source_id is provided, uses it; otherwise creates a new one.
        Returns the chat_history_id (visual_source_id).
        """
        # If visual_source_id not provided, create one
        if visual_source_id is None:
            # Always create visual_source to satisfy foreign key constraint
            # The database requires chat_history_id to reference a valid visual_source_id
            from sqlyst_app.crud.visual_source import insert_visual_source
            visual_source = insert_visual_source(db, "CHAT", None)
            # Commit visual_source to ensure it exists before creating chat_history with foreign key reference
            db.commit()
            chat_history_id = visual_source.visual_source_id
        else:
            chat_history_id = visual_source_id

        # Save chat history
        crud.create_chat_history(ChatHistoriesRequest(
            session_id=session_id,
            user_message=query,
            model_response=model_response
        ), db, chat_history_id=chat_history_id)

        return chat_history_id


    def chat(self, query: str, session_id: UUID, db: Session, user_id: UUID, user_name: str = None) -> ChatResponse:
        """Process chat query and return response with follow-up questions and visualization data."""

        session = session_crud.get_session_by_id_and_user_id(session_id, user_id, db)
        if not session:
            raise NotFoundException(f"Session not found for user: {user_id}")

        history_saved = False
        try:
            # Get chat history context
            context = crud.get_context_chat_histories_by_session_id(session_id, db)

            # Convert to BaseMessage format
            chat_messages = self._convert_chat_history_to_messages(context)

            # Add current query
            current_message = HumanMessage(content=query)

            # Invoke agent with proper message format
            response = get_agent().invoke({
                "messages": [current_message],
                "chat_history": chat_messages,
                "user_name": user_name
            })


            print("DEBUG: response",response)

            if not response or not response.get("messages"):
                raise ValueError("Agent returned empty response")

            # Extract response data safely
            messages = response.get("messages", [])

            print("DEBUG: messages",messages)

            explanation = messages[-1].content if messages else "No response generated"

            print("DEBUG: explanation before extract_json_from_response",explanation)
            # Extract visualization data and follow-up questions from the raw response
            viz_json = extract_json_from_response(explanation)

            print("DEBUG: viz_json",viz_json)

            # Extract follow-up questions before cleaning - check for FOLLOW_UP_QUESTIONS section
            has_follow_ups = (
                "FOLLOW_UP_QUESTIONS" in explanation or
                "\uD6C4\uC18D \uC9C8\uBB38" in explanation or  # "Follow-up question" (Korean)
                re.search(r'#{1,6}\s*FOLLOW_UP_QUESTIONS', explanation, re.IGNORECASE)
            )
            follow_up_questions = extract_follow_up_questions_from_response(explanation) if has_follow_ups else []

            # Clean explanation: keep INSIGHTS if exists, otherwise keep plain text (removes FOLLOW_UP_QUESTIONS)
            explanation = keep_only_insights(explanation)

            # Store visual in item_container if visualization data exists
            item_id = None

            # Always create visual_source to satisfy foreign key constraint
            # The database requires chat_history_id to reference a valid visual_source_id
            from sqlyst_app.crud.visual_source import insert_visual_source
            visual_source = insert_visual_source(db, "CHAT", None)
            # Commit visual_source to ensure it exists before creating chat_history with foreign key reference
            db.commit()

            # Use visual_source_id as chat_history_id (same as forecast_id = visual_source_id)
            chat_history_id = visual_source.visual_source_id

            if viz_json:
                from sqlyst_app.crud.visuals import insert_visual_from_forecast_and_chat_response

                # Prepare visual data
                viz_json["is_stored"] = False
                chart_type = viz_json.get("chart_type", "line")
                if chart_type not in ["bar", "line", "pie", "table"]:
                    chart_type = "line"
                # Add "chat_" prefix to chart_type
                if not chart_type.startswith("chat_"):
                    chart_type = f"chat_{chart_type}"
                viz_json["chart_type"] = chart_type

                # Insert visual using the same function as forecast
                visual = insert_visual_from_forecast_and_chat_response(
                    db, viz_json, visual_source.visual_source_id, user_id
                )
                item_id = visual.item_id
                viz_json["item_id"] = str(item_id)

            # Follow-up questions already extracted above

            # Create clean response for storage
            clean_response = {
                "explanation": explanation,
                "sql_query": response.get("sql_query", ""),
            }
            if item_id:
                clean_response["visualization_item_id"] = str(item_id)

            # Save chat history with chat_history_id (from visual_source_id if visual exists)
            self._save_chat_history(query, session_id, clean_response, db, chat_history_id)
            history_saved = True

            chat_response = ChatResponse(
                explanation=explanation,
                follow_up_questions=follow_up_questions,
                visualization_data=viz_json
            )

            return chat_response

        except Exception as e:
            # Handle errors gracefully
            print(f"ERROR in service: {str(e)}")

            # Create error response message
            error_message = f"There is an error processing your request. Please try again."

            # Create error response for storage
            error_response = {
                "explanation": error_message,
                "sql_query": "",
                "error": str(e)
            }

            # Save error to chat history to maintain conversation flow
            try:
                self._save_chat_history(query, session_id, error_response, db)
                history_saved = True
            except Exception as save_error:
                # If saving to history fails, log but don't fail the request
                print(f"ERROR saving to chat history: {str(save_error)}")

            return ChatResponse(
                explanation=error_message,
                follow_up_questions=[],
                visualization_data=None
            )
        finally:
            # Save timeout error if history wasn't saved (e.g., timeout interrupted execution)
            if not history_saved:
                try:
                    timeout_response = {
                        "explanation": "Request timeout. The AI is taking longer than expected to process your query. Please try again with a simpler question.",
                        "sql_query": "",
                        "error": "TimeoutError"
                    }
                    self._save_chat_history(query, session_id, timeout_response, db)
                except Exception as save_error:
                    print(f"ERROR saving timeout to chat history: {str(save_error)}")
\end{lstlisting}

% ---------------------------------------------------------------------
\subsection{LLM Factory}
\label{sub:llm-factory}

\code{llm\_factory.py} centralizes construction of the language,
Text2SQL, and embedding model clients used across the agent pipeline,
lazily loading the local Ollama models and exposing thin
generation/embedding helpers around the vLLM-served Text2SQL and
embedding endpoints.

\begin{lstlisting}[style=python, caption={Model client factory (\code{llm\_factory.py})}, label={lst:llm-factory}]
import os

import openai
from langchain_ollama import OllamaLLM


def _ollama_base_url() -> str:
    """Resolve the Ollama base URL, defaulting to the host gateway."""
    return os.environ.get("OLLAMA_BASE_URL", "http://localhost:11435")

# Global model instances for lazy loading (only for Ollama models)
_translator_model = None
_reasoning_model = None
def translatorModel():
    """Lazy load the translator model."""
    global _translator_model
    if _translator_model is None:
        _translator_model = OllamaLLM(
            model="timHan/llama3korean8B4QKM:latest",
            base_url=_ollama_base_url(),
            timeout=20,
        )
    return _translator_model

# Text2SQL model configuration
TEXT2SQL_VLLM_API_BASE = os.environ.get("VLLM_API_BASE", "http://localhost:8007/v1")
TEXT2SQL_MODEL_NAME = "Snowflake/Arctic-Text2SQL-R1-7B"

# Embedding model configuration
EMBED_VLLM_API_BASE = os.environ.get("EMBED_VLLM_API_BASE", "http://localhost:8008/v1")
EMBED_MODEL_NAME = "Snowflake/snowflake-arctic-embed-m"

def text2SqlModel():
    """Get the Text2SQL model client."""
    return openai.OpenAI(
        api_key="sk-not-needed",
        base_url=TEXT2SQL_VLLM_API_BASE
    )


def text2sql_generate(
    prompt: str,
    max_tokens: int = 400,
    temperature: float = 0.0
) -> str:
    """
    Convenience method to generate SQL using the Text2SQL model (completion model).

    Args:
        prompt: Complete prompt string (can include instructions, schema, user question, etc.)
        max_tokens: Maximum tokens to generate (default: 200)
        temperature: Sampling temperature (default: 0.0 for deterministic SQL)

    Returns:
        Generated SQL query as string
    """
    client = text2SqlModel()

    response = client.completions.create(
        model=TEXT2SQL_MODEL_NAME,
        prompt=prompt,
        max_tokens=max_tokens,
        temperature=temperature
    )
    return response.choices[0].text

def reasoningModel() -> OllamaLLM:
    """Lazy load the reasoning model."""
    global _reasoning_model
    if _reasoning_model is None:
        print("Loading reasoning model")
        model_name = "llama3.1:8b"
        _reasoning_model = OllamaLLM(
            model=model_name,
            base_url=_ollama_base_url(),
            timeout=30,
            num_gpu=-1
        )
        print("Reasoning model loaded successfully")

    return _reasoning_model

def embedModel():
    """Get the embedding model client."""
    return openai.OpenAI(
        api_key="sk-not-needed",
        base_url=EMBED_VLLM_API_BASE
    )


def embed_generate(
    texts: list[str] | str
) -> list[list[float]]:
    """
    Convenience method to generate embeddings using the embedding model.

    Args:
        texts: Single text string or list of text strings to embed

    Returns:
        List of embedding vectors (each as a list of floats)
    """
    client = embedModel()

    # Ensure texts is a list
    if isinstance(texts, str):
        texts = [texts]

    response = client.embeddings.create(
        model=EMBED_MODEL_NAME,
        input=texts
    )

    # Extract embeddings from response
    embeddings = [item.embedding for item in response.data]
    return embeddings
\end{lstlisting}
